\section{The boundary}
\label{sec:boundary}

\Cref{fig:ladder} contains an uncomfortable fact. The mechanism, built from two
measured parameters, moves the \ProjectionEndYear{} world result by
\LadderMainstreamGain{} billion. A single unmeasured scenario knob --- how fast
the fertility environment keeps falling --- moves it by \LadderPressureLoss{}
billion the other way. Reported as a preferred value, that knob would be doing
most of the work while carrying none of the evidence, and the paper would amount
to an assertion dressed as a projection.

There is a better use for it, available because of a structural property of the
model.

\subsection{Selection and the environment are commensurable}
\label{sec:separability}

The environment multiplier $\envmult_t[\ell]$ in \Cref{eq:births-by-mother}
multiplies every type equally. It therefore cancels from the relative weights
with which types contribute to the next generation: whether conditions are
generous or harsh, the ratio of a high-propensity woman's births to a
low-propensity woman's births is unchanged, so the composition follows almost the
same path either way.

The cancellation is exact in the birth weights and very nearly exact in the
outcome. Running the same selection setting under two environments that end
\SeparabilityPopulationGap{} billion people apart gives selection multipliers of
\SeparabilityCalmEffect{} and \SeparabilityStressedEffect{}, a difference of
\SeparabilityGap{}. The residue is a real second-order effect rather than
rounding: a harsher environment produces a different age structure, so the
childbearing-age women over whom the propensity is averaged are drawn
differently, and a fixed number of migrants mixes into a smaller population with
proportionally more force.

That near-separability is a modelling assumption --- \Cref{sec:limitations} takes
it seriously as one --- but it is what makes the two forces directly comparable.
Selection multiplies fertility by $S$; a uniform environmental decline of rate
$r$ per decade over $d$ decades multiplies it by $(1-r)^d$. The two cancel when
\begin{equation}
  S \cdot (1-r)^{d} = 1
  \qquad\Longleftrightarrow\qquad
  r^{\star} = 1 - S^{-1/d}.
  \label{eq:break-even}
\end{equation}

\subsection{Two thresholds, because there are two questions}

\Cref{eq:break-even} cancels the fertility \emph{rate} in the final year. At the
central calibration, measured mainstream selection reaches
$S = \MainstreamSelectionMultiplier{}$ in the final fertility-rate year
(\FinalRateYear{}), giving \BoundaryDeterministic{} per decade of additional
decline beginning in \PressureFromYear{}. Country selection effects are
aggregated using births from the no-selection counterfactual, so a scenario
cannot choose the weights by which it is judged.

That is not the same as cancelling selection's effect on the population, and the
difference is not small. Bringing the rate back to parity in \FinalRateYear{}
leaves every extra birth selection produced in the preceding century, and those
people and their descendants are alive in \ProjectionEndYear{}. Solving directly
for the decline that returns the \ProjectionEndYear{} world population to the
no-selection benchmark of \LadderBenchmark{} billion gives
\BoundaryPopulation{} per decade --- sixty per cent larger. Both are reported
because both get asked, and quoting only the first would understate what it takes
to undo the mechanism.

\begin{figure}[t]
\centering
\includegraphics[width=0.86\textwidth]{fig3-boundary.pdf}
\caption{The additional uniform decline in the fertility environment, per decade
after \PressureFromYear{}, that cancels mainstream selection's effect on the
fertility rate by \ProjectionEndYear{}. Above a curve the environment wins;
below it selection does. The two axes are the measured parameters: the dispersion
of completed family size (three curves, spanning the observed low-fertility
country range) and the parent--child correlation in it. Shaded bands are the 5th
to 95th percentile across \PairedMigrationPaths{} paired stochastic migration
paths, and are narrow enough to be nearly invisible, which is itself the result
reported in \Cref{sec:migration-robustness}. The dashed line marks the
illustrative \PressureStressTest{} stress test of \Cref{fig:ladder}. Cancelling
the effect on the 2150 \emph{population} rather than on the rate requires
\BoundaryPopulation{} at the central setting, above the top of this figure's
central curve.}
\label{fig:boundary}
\end{figure}

\Cref{fig:boundary} shows how the rate threshold moves across the empirical range
of both measured parameters, over a \BoundaryGridCells{}-cell grid. It runs from
\BoundaryFloor{} per decade at the low corner --- weak dispersion and weak
transmission --- to \BoundaryCeiling{} per decade at the high corner. The
threshold is close to linear in the transmission coefficient and roughly
quadratic in the dispersion, as \Cref{eq:response} implies.

Read as a statement about the world, this says something more useful than a
projected total. If the conditions that determine how many children a given
person has continue to deteriorate at less than roughly one and a half per cent
per decade after \PressureFromYear{}, the compositional force wins on the
fertility rate; below about two and a half per cent it still wins on the
population at the horizon. Above those, the environment wins. Nobody can
currently say which, and this paper does not try to. What it does is put a number
on the question, so that a claim about the future of fertility can be stated as a
rate that could in principle be observed.

For scale: \PressureStressTest{} per decade, the illustrative stress test, is
well above both thresholds and is far from a neutral choice. Sustained for a
century it removes a third of the fertility environment before its floor binds.
It is reported because it was the project's original central scenario and a
stress test should be visible, not because we think it likely.

\subsection{The boundary survives stochastic migration}
\label{sec:migration-robustness}

A threshold computed on one migration path would be worth little if migration
could move it. It is checked against \PairedMigrationPaths{} stochastic paths
drawn from the published Bayesian migration trajectories \citep{azose2015}, with
one design choice that matters: the selection run and its no-selection
counterfactual share each source path and each post-2100 innovation, so the
comparison is paired and a migration draw cannot masquerade as a selection
effect. Every draw-year is balanced to exactly zero world net migration.

The central rate threshold is \BoundaryPaired{} per decade with a 90\% range of
\BoundaryPairedLow{} to \BoundaryPairedHigh{}. Migration moves it by about one
part in seventy-five, against a factor of thirteen between the corners of the
measured-parameter grid. That is what one should expect --- migration sums to
zero globally by construction and can affect a world fertility aggregate only by
relocating people between countries with different schedules --- but expecting it
is not the same as having checked it.

\subsection{What the boundary is not}

\caveatbox{Reading \Cref{fig:boundary} correctly}{The vertical axis is
\emph{additional} decline, on top of the fertility fall already in the observed
data and already carried by the benchmark path. It is not a forecast of economic
conditions, and this paper makes no claim about what moves the fertility
environment or how fast. The threshold is also specific to the terminal year:
because both forces compound, a decline beginning later or ending earlier trades
against the same total. And the uniformity is deliberate --- a decline falling
unequally across dispositions would interact with selection rather than merely
oppose it, which is a richer model than the evidence supports and which would
destroy the single threshold that makes this figure readable. \Cref{sec:limitations}
treats that as the most valuable available extension.}
